% Source: 01 Properties_of_Matter 26(1).docx; formatting in layout.tex.
\HOne{The Properties of Substances}

\textbf{Ozone }is a gas that plays an important role in the stratosphere, where it absorbs much of the Sun's harmful ultraviolet radiation. At ground level, however, ozone is an air pollutant. It forms from other pollutants in strong sunlight and can irritate the lungs.


\textbf{Penicillin }is a white powder. It is used as an antibiotic. It is a very powerful drug and one reason for the increase in life expectancy in the 20th century.


\textbf{Arsenic }is toxic in very small amounts. It has been used in the past to get rid of unwanted heirs or heiresses to the throne.


\textbf{Ethanol} still is the most widely used substance which causes addiction.


\textbf{Copper} is used in electrical wires. Without copper not a single computer would work.


Although there are many different substances, each substance has a \textbf{characteristic combination of properties}. Some measured properties depend on external conditions, especially temperature and pressure. For example, more sugar dissolves in hot water than in cold water, and the density of a gas increases when the pressure increases.


\HTwo{How Scientists Measure Properties}

Some properties can be described without a measurement. For example, a substance may be colourless, shiny, solid, or brittle. Other properties are expressed using measurements. A measurement consists of a numerical value and a unit.


\begin{StandardTable}{
  colspec={X[l,m] X[l,m] X[l,m]}
}

{Term} & {Meaning} & {Example} \\ 

{Physical quantity} & {what is being measured} & {mass} \\ 

{Numerical value} & {the measured number} & {2.5} \\ 

{Unit} & {the agreed reference} & {g}

\end{StandardTable}


\begin{InfoBox}
\BoxHeading{Measurement}
A measurement is a numerical value times a unit: measurement = numerical value $\cdot$ unit \newline \textit{Example: m = 2.5 $\cdot$ g}
\end{InfoBox}


Measurements can only be compared reliably when the units are clearly defined. Scientists, therefore, use the International System of Units, abbreviated SI. It provides a common language for measurements in science, technology, and everyday life.


\begin{StandardTable}{
  colspec={X[l,m] X[l,m] X[l,m]}
}

{Quantity} & {SI Unit} & {Symbol} \\ 

{length (l)} & {metre} & {m} \\ 

{mass (m)} & {kilogram} & {kg} \\ 

{time (t)} & {second} & {s} \\ 

{electric current} & {ampere} & {A} \\ 

{temperature (T)} & {kelvin} & {K} \\ 

{amount of substance (n)} & {mole} & {mol} \\ 

{luminous intensity} & {candela} & {cd}

\end{StandardTable}

\newpage

In this chapter, we will mainly use units for mass, length, temperature, and time. The mole will become important in a later chapter.


The SI system contains seven base quantities with seven corresponding base units. Other physical quantities are described using derived units formed by multiplying or dividing the base units. The unit can be directly \textbf{derived from the formula}.


\textbf{Area: }A = l $\cdot$ l ; [A] = m $\cdot$ m = m$^2$


\textbf{Volume: }V = l $\cdot$ l $\cdot$ l ; [V] = m $\cdot$ m $\cdot$ m = m$^{3}$


\textbf{Density: }$\rho$ = m / V ; [$\rho$] = kg / m$^3$


In the laboratory, units such as grams, millilitres, cubic centimetres, and degrees Celsius are often more convenient than the corresponding SI units.


\begin{StandardTable}{
  colspec={X[l,m] X[l,m] X[l,m]}
}

{Quantity} & {SI unit} & {Common laboratory units} \\ 

{Mass} & {kg} & {g, mg} \\ 

{Volume} & {m$^3$} & {L = dm$^3$, mL = cm$^3$} \\ 

{Temperature} & {K} & {$^{\circ}$C} \\ 

{Time} & {s} & {min}

\end{StandardTable}


The same unit may be used with a \textbf{prefix} when a value is very large or very small. A prefix changes the unit by a power of ten.


\begin{StandardTable}{
  colspec={X[l,m] X[l,m] X[l,m] X[l,m]}
}

{Prefix} & {Symbol} & {Factor} & {Power of ten} \\ 

{kilo} & {k} & {1'000} & {10$^3$} \\ 

{centi} & {c} & {0.01} & {10$^{-2}$} \\ 

{milli} & {m} & {0.001} & {10$^{-3}$} \\ 

{micro} & {$\mu$} & {0.000 001} & {10$^{-6}$}

\end{StandardTable}


\begin{InfoBox}
\BoxHeading{Prefix and Scientific Notation}
Very large and very small values can be written compactly using SI prefixes, scientific notation, or a combination of both.\newline \textbf{Prefix notation}: 0.000 001 m = 1 $\mu$m\newline \textbf{Scientific notation}: 40 000 000 m = 4  10$^7$ m\newline \textbf{Combination}: 40 000 000 m = 4  10$^4$ km
\end{InfoBox}

\newpage

\HTwo{States of Matter}

\begin{StandardTable}{
  colspec={X[l,m] X[l,m] X[l,m] X[l,m] X[l,m]},
  row{2-4}={ht=18mm}
}

 & volume & density & shape & fluidity \\

\textbf{solid} & fixed volume & & & \\

\textbf{liquid} & & & no defined shape & \\

\textbf{gaseous} & & low & & flows easily

\end{StandardTable}


\HFour{The Particle Theory of Matter}

According to the particle theory of matter all substances are made up of particles. The particles are extremely small, so we cannot see them. In a certain substance, for example water, all particles are the same. Particles of different substances are different from each other.


\begin{minipage}[t]{0.60\linewidth}
  \vspace{0pt}
    When a substance is in its \textbf{solid state}, the particles touch each other and are in a regular pattern. The particles are very strongly attracted to each other. They don't move around but vibrate on the spot.

    The letter used for the solid state is \textbf{(s)}.
\end{minipage}\hfill
\begin{minipage}[t]{0.35\linewidth}
  \vspace{0pt}
  \begin{DrawingBox}{\linewidth}{36mm}
  \end{DrawingBox}
\end{minipage}


\begin{minipage}[t]{0.60\linewidth}
  \vspace{0pt}
    When a substance is in its \textbf{liquid state}, the particles are free to move around, but they are still in contact. The particles are strongly attracted to each other, but there are no regular patterns. Under the influence of gravity, the particles slide off each other.
    
    The letter used for the liquid state is \textbf{(l)}.
\end{minipage}\hfill
\begin{minipage}[t]{0.35\linewidth}
  \vspace{0pt}
  \begin{DrawingBox}{\linewidth}{36mm}
  \end{DrawingBox}
\end{minipage}


\begin{minipage}[t]{0.60\linewidth}
  \vspace{0pt}
    When a substance is in its \textbf{gaseous state}, the particles don't touch each other. The forces of attraction between the particles are very weak. The particles move fast in all directions. There is no regular pattern.

    The letter used for the gaseous state is \textbf{(g)}.
\end{minipage}\hfill
\begin{minipage}[t]{0.35\linewidth}
  \vspace{0pt}
  \begin{DrawingBox}{\linewidth}{36mm}
  \end{DrawingBox}
\end{minipage}


% QR-Code mit Link zur Simulation
\begin{minipage}[c]{0.13\linewidth}
  \includegraphics[width=\linewidth]{images/image18.png}
\end{minipage}\hfill
\begin{minipage}[c]{0.83\linewidth}
  \href{https://phet.colorado.edu/sims/html/states-of-matter-basics/latest/states-of-matter-basics_all.html?locale=en}
  {\textcolor{DeepBlue}{\underline{Here}}},
  a simulation can be found to visualize the states of matter.
\end{minipage}


\HTwo{Melting and Boiling Points}

The \textbf{melting point} (mp) is the temperature at which a solid and its liquid state coexist while the substance melts. For a pure substance at a given pressure, melting occurs at a \textbf{characteristic temperature}. The melting point can therefore help identify an unknown substance.


The reverse process is \textbf{freezing}. At the same pressure, a pure substance freezes at the same temperature at which it melts. Pressure usually has only a small effect on the melting point of most substances.


\Exercise[ssdfsd]{sdvfsdf}

\Exercise[bar-chart]{The bar chart shows the melting points of eight metals.}

\MetalBarChart


\begin{ExerciseParts}

\item Name the metal with the lowest melting point.

\ResponseGap[12pt]

\item List the metals in order of increasing melting point, beginning with the lowest.

\ResponseGap

\item An unknown metal has a measured melting point of 1050 $^{\circ}$C. Which metal in the chart is the most likely match? Explain your choice.

\ResponseGap[20pt]

\end{ExerciseParts}

\begin{wrapfigure}{r}{0.34\linewidth}
  \centering
  % Replaces images/image20.jpg (original retained)
\resizebox{\linewidth}{!}{\GasBubblesInLiquidFigure}
\end{wrapfigure}

The \textbf{boiling point} (bp) is the temperature at which a liquid boils throughout the liquid. During boiling, bubbles containing the substance in the gaseous state form inside the liquid, rise to the surface, and escape.


\textbf{Boiling} begins when the vapour pressure of the liquid equals the surrounding pressure. The boiling point, therefore, depends strongly on pressure. At standard atmospheric pressure (1013 hPa), pure water boils at 100 $^{\circ}$C. At higher altitudes, atmospheric pressure is lower, so water boils at a lower temperature.


The reverse process is \textbf{condensation}. At the same pressure, a pure substance condensates at the same temperature at which it boils.


\textbf{Evaporation} is the change from a liquid to a gas at the surface of the liquid. It can occur at temperatures below the boiling point. In contrast, boiling occurs throughout the liquid.


A substance can also change directly from a solid to a gas. This process is called \textbf{sublimation}. The reverse change, from a gas directly to a solid, is called \textbf{resublimation}.


\Exercise[mp-table]{\textbf{The table shows the melting and boiling points of several substances at standard atmospheric pressure.}}

\begin{StandardTable}{
  colspec={X[l,m] X[c,m] X[c,m]}
}

substance & melting point / $^{\circ}$C & boiling point / $^{\circ}$C \\

{sodium} & {98} & {883} \\ 

{radon} & {- 71} & {- 62} \\ 

{ethanol} & {- 114} & {78} \\ 

{cobalt} & {1495} & {2927} \\ 

{nitrogen} & {- 210} & {- 196} \\ 

{propane} & {- 187} & {- 42} \\ 

{ethanoic acid} & {17} & {118}

\end{StandardTable}


\begin{ExerciseParts}

\item Classify each substance as solid, liquid, or gaseous at 25 $^{\circ}$C and standard atmospheric pressure.

\ResponseGap[20pt]

\item Which substance has the smallest liquid temperature range?

\ResponseGap[20pt]

\item Which substances are gaseous at -50 $^{\circ}$C?

\ResponseGap[20pt]

\item Which substance has the lowest melting point?

\ResponseGap[20pt]

\item Which substance is liquid at 2500 $^{\circ}$C?

\ResponseGap[20pt]

\end{ExerciseParts}

\Exercise{Draw the changes of state in the particle model in the white boxes.}

\begin{center}% Replaces images/image21.png (original retained)
\resizebox{0.70\linewidth}{!}{\StateChangesWorksheetFigure}\end{center}


\HTwo{Solubility}

When sugar is added to tea, it appears to disappear. However, the sugar is still present: it has \textbf{dissolved}. In this example, the tea acts as the \textbf{solvent}, the sugar is the \textbf{solute}, and the sweetened tea is a \textbf{solution}.


Only a limited amount of a solute can dissolve under given conditions. If more sugar is added than can dissolve, some remains undissolved. The solution is then \textbf{saturated}.


A solution is a homogeneous mixture consisting of a solvent and one or more solutes. A solution in which water is the solvent is called an aqueous solution. A substance is \textbf{soluble} in a particular solvent if a noticeable amount dissolves; otherwise, it is described as \textbf{insoluble} or only slightly soluble.


A saturated solution contains the maximum amount of dissolved solute possible at a specified temperature. Additional solute remains undissolved.


\begin{InfoBox}
\BoxHeading{Solubility}
The maximal amount of a substance in grams (or ml at standard pressure if the substance is a gas) which can be dissolved in 100 ml of the solvent.
\end{InfoBox}


Solubility depends on the substances involved and often on temperature. The values below compare the solubility of several substances in water at 20 $^{\circ}$C. 


\begin{center}
\begin{tabular}{@{}l@{\hspace{25mm}}r@{}}
  table salt (sodium chloride) & 35.9 \\
  limestone (calcium carbonate) & 0.002 \\
  ethanol & any amount \\
  bromine & 3.8 \\
  ammonia & 74'000 \\
  oxygen & 3.2
\end{tabular}
\end{center}


\Exercise[sol-assign]{\textbf{Match each term with its correct definition.}}

\begin{tblr}{
  width=\linewidth,
  colspec={
    Q[c,m,wd=9mm] X[1.4,l,m]
    Q[wd=9mm]
    Q[c,m,wd=9mm] X[2.2,l,m]
  },
  colsep=4pt,
  rowsep=3pt,
  hline{1-Z}={1-2,4-5}{0.8pt,TableGrid},
  vline{1-3}={0.8pt,TableGrid},
  vline{4-6}={0.8pt,TableGrid}
}

A & solute     & & 1 & The process of adding a solid to a liquid so that you can no longer see separate pieces of the solid. \\
B & solvent    & & 2 & A substance that is dissolved in a liquid. \\
C & solution   & & 3 & A substance is \ldots{} in a liquid if it dissolves in the liquid. \\
D & dissolving & & 4 & The mass of substance that dissolves in 100 ml of water. \\
E & soluble    & & 5 & A liquid in which another substance is dissolved. \\
F & insoluble  & & 6 & A substance is \ldots{} in a liquid if it does not dissolve in the liquid. \\
G & solubility & & 7 & A mixture of a liquid and a dissolved substance.

\end{tblr}

\newpage

\Exercise{Draw the dissolving process in the particle model}

\ResponseGap[100pt]

\Exercise[sol-graph]{\textbf{Use the solubility graph to answer the following questions.}}

\begin{minipage}[t]{0.55\linewidth}
  \vspace{0pt}

  \begin{ExerciseParts}

    \item What difference can you observe
    between the solubility of gases (\ce{HCl}, \ce{SO2}, \ce{NH3},…) and solids.

    \ResponseGap

    \item Compare how the solubility of \ce{KNO3} and \ce{NaCl}
    changes with temperature.

    \ResponseGap

	\item According to the graph, how many grams of \ce{NaNO3}
    dissolve in 100 g of water at room temperature?
    Explain why the graph alone does not directly give
    the mass required to prepare 250 mL of saturated solution.

    \ResponseGap[60pt]

  \end{ExerciseParts}

\end{minipage}\hfill
\begin{minipage}[t]{0.42\linewidth}
  \vspace{0pt}
  %\includegraphics[width=\linewidth]{images/image22.png}
  \SolubilityCurvesFigure
\end{minipage}

\Exercise[sol-sat]{\textbf{A saturated \ce{KCl} solution at 80 $^{\circ}$C is cooled to room temperature. Describe what happens and represent the process using the particle model.}}

\ResponseGap[60pt]

\newpage
\HTwo{Density}

An iron dumbbell can have the same mass as a much larger wooden object. The reason is density. Density describes how much mass is contained in a given volume. The Greek letter \textbf{$\rho$} (rho) is used as the symbol for \textbf{density}.

\[\rho=\frac{m}{V}\]


The SI unit of density is \textbf{kg/m$^3$}. In the laboratory, \textbf{g/cm$^3$} and \textbf{g/mL} are often more convenient. Since 1 mL = 1 cm$^3$, these two units have the same numerical value.


Density depends on temperature, pressure, and the state of matter. The effect of pressure is especially important for gases. Density values should, therefore, be reported together with the relevant conditions. Under specified conditions, density is a characteristic property that can help identify a substance.


Examples at 20 $^{\circ}$C and standard atmospheric pressure:


\begin{center}
\begin{tabular}{@{}l@{\hspace{20mm}}r@{}}
  Gold           & 19.3 $\cdot 10^3$ kg/m$^3$ \\
  Silver         & 10.5 $\cdot 10^3$ kg/m$^3$ \\
  Ethanol        & 0.789 $\cdot 10^3$ kg/m$^3$ \\
  Mercury        & 13.5 $\cdot 10^3$ kg/m$^3$ \\
  Oxygen         & 1.33 kg/m$^3$ \\
  Carbon Dioxide & 1.84 kg/m$^3$
\end{tabular}
\end{center}


\Exercise[dens-table]{\textbf{Complete the missing values in the table.}}

\begin{StandardTable}{
  colspec={X[l,m] X[c,m] X[c,m] X[c,m]}
}

substance & mass / g & volume / cm$^3$ & density / (g/cm$^3$) \\

balsa wood & 16.0 & 80.0 & 0.200 \\

oak wood & 42.0 & 60.0 & \\

copper & 17.8 & & 8.90 \\

aluminium & & 4.00 & 2.70

\end{StandardTable}


\Exercise[dens-cal]{\textbf{A 50.00 mL sample of a liquid has a mass of 62.37 g. Calculate its density.}}

\ResponseGap[50pt]

\Exercise[dens-vol]{\textbf{A gold crown and a silver crown each have a mass of 10.00 kg. Calculate the volume of water displaced when each crown is completely submerged.}}

\ResponseGap[80pt]

\newpage

\Exercise[dens-unit]{\textbf{Express the density of gold in the following units:}}

\begin{ExerciseParts}

\item g/cm$^3$
\ResponseGap

\item kg/L
\ResponseGap

\item mg/$\mu$L
\ResponseGap

\end{ExerciseParts}

\Exercise[dens-beam]{\textbf{A piece of silver is placed on one side of a beam balance. The same volume of mercury is placed on the other side. }\textbf{Predict what will happen and support your answer with a calculation or a labelled drawing.}}

\ResponseGap[80pt]

\Exercise[dens-gly]{\textbf{An experiment requires 38.75 g of glycerol. No balance is available, so you must measure the glycerol using a graduated cylinder.}\textbf{ }\textbf{Calculate the required volume in millilitres. ($\rho$(glycerol) = 1.261 }\textbf{$\cdot$}\textbf{ 10$^3$ kg/m$^3$)}}

\ResponseGap[80pt]

\Exercise[dens-exp]{A student wants to identify an unknown metal.
The metal has a mass of 53.4 g. When placed in a
graduated cylinder, the water level rises from
35.0 mL to 41.0 mL.}

\begin{ExerciseParts}

\item Calculate the density of the metal and identify
it using the densities given earlier in this section.

\ResponseGap[45pt]

\item During the experiment, a small air bubble
remains attached to the metal.
Explain how this affects the calculated density.

\ResponseGap[45pt]

\end{ExerciseParts}

\newpage
\HTwo{Other Properties}

Not all characteristic properties are numerical. Chemists also describe substances using observable or testable properties such as colour, odour, electrical conductivity, thermal conductivity, brittleness, and malleability.


\Exercise[other-assign]{\textbf{Match each property with the corresponding picture.}}


\begin{tblr}{
  width=\linewidth,
  colspec={
    Q[c,m,wd=9mm] X[1.7,c,m]
    Q[wd=40mm]
    Q[c,m,wd=9mm] X[1.7,c,m]
  },
  colsep=4pt,
  rowsep=3pt,
  % Zwei getrennte Tabellenbereiche
  hline{1-Z}={1-2,4-5}{0.8pt,TableGrid},
  vline{1-3}={0.8pt,TableGrid},
  vline{4-6}={0.8pt,TableGrid},
  % Genügend Platz für die Bilder
  row{1-9}={ht=12mm}
}

{A} & {colour} & {} & {1} & \raisebox{-0.45\height}{\includegraphics[height=11mm]{images/image24.png}} \\ 

{B} & {ductility/ malleability} & {} & {2} & \raisebox{-0.45\height}{\includegraphics[height=11mm]{images/image25.jpeg}} \\ 

{C} & {thermal conductivity} & {} & {3} & \raisebox{-0.45\height}{\includegraphics[height=11mm]{images/image26.png}} \\ 

{D} & {solubility} & {} & {4} & \raisebox{-0.45\height}{\includegraphics[height=11mm]{images/image27.jpeg}} \\ 

{E} & {odor} & {} & {5} & \raisebox{-0.45\height}{\includegraphics[height=11mm]{images/image28.png}} \\ 

{F} & {electrical conductivity} & {} & {6} & \raisebox{-0.45\height}{\includegraphics[height=11mm]{images/image29.png}} \\ 

{G} & {brittleness} & {} & {7} & \raisebox{-0.45\height}{\includegraphics[height=11mm]{images/image30.jpeg}} \\ 

{H} & {density} & {} & {8} & \raisebox{-0.45\height}{\includegraphics[height=11mm]{images/image31.jpeg}} \\ 

{I} & {state of matter at room temperature} & {} & {9} & \raisebox{-0.45\height}{\includegraphics[height=11mm]{images/image32.png}}

\end{tblr}

\par\smallskip


All the properties discussed in this chapter are \textbf{physical properties}. Substances also have \textbf{chemical properties}, which describe how they react with other substances. These properties will be discussed in a later chapter.

